\chapter{Integrity}
\label{ch:integrity}

\section{The failure mode}

Every chapter of Part III has a way of going wrong. This is the one that is
common to all of them, and it is the way transitions of this kind actually die.

The mechanism is not mysterious. A transition takes assets that were previously
unpriced --- land, buildings, licences, spectrum, ports, mineral rights, state
enterprises --- and prices them, over a period of a few years, in a country whose
courts are weak, whose press is new, whose officials are underpaid, and whose
rules are being rewritten. During that window, a small number of people are in a
position to decide who gets what, and they know the window will close. Everything
about the situation rewards taking as much as possible as fast as possible.

The results are on the record. Russia converted a planned economy into a private
one in five years and produced, by the end of it, an ownership structure so
concentrated and so illegitimate that the political system reorganised itself
around protecting it. Ukraine spent twenty-five years with a formally democratic
system captured by a handful of industrial groups \citep{aslund2013}. Romania's privatisations
produced a decade of scandal. Against those, Poland, Estonia and Slovenia did not
avoid corruption --- nobody does --- but they kept it below the level at which it
determines who governs, and they did so with a set of instruments that are
well-documented and unglamorous.

Cuba's version of the risk is specific and named in
Chapter~\ref{ch:premises}: the incumbents hold the assets, the accounts and the
information, and they are the people who would be administering the transition. A
liberalisation conducted before the checks exist does not distribute the Cuban
economy. It hands it to the people currently holding it, with legal title
attached, and that outcome is effectively permanent.

\section{Two different problems}

An important distinction, because Cuba's public discussion and its criminal law
run them together and the instruments for each are opposite.

\emph{Survival corruption.} In an economy of chronic shortage, ordinary Cubans
have spent thirty years obtaining what they need through informal channels ---
\emph{resolver}, \emph{luchar}, the state warehouse worker who diverts a case of
cooking oil, the doctor given a gift for an appointment, the inspector who
overlooks a thing. This is pervasive, it is how the country has fed itself, and
it is a symptom of scarcity and of wages that do not cover subsistence. It is not
addressed by enforcement; forty years of enforcement have not addressed it. It is
addressed by supply and by wages, which is to say by Chapters~\ref{ch:land},
\ref{ch:energy} and \ref{ch:welfare}. Treating it as the same offence as the next
category is both unjust and a distraction, and it has the further cost of making
the whole population complicit and therefore uninterested in anti-corruption as
an idea.

\emph{State capture.} A small number of decisions \citep{hellman1998} --- who gets the port
concession, who gets the currency at the official rate, who buys the hotel chain,
whose company wins the solar tender --- determine the ownership of the country.
This is the category that decides whether Cuba in 2045 is Estonia or Ukraine, and
it is entirely about a few hundred people and a few thousand decisions.

Everything in this chapter is aimed at the second. The first is a welfare
problem wearing a criminal-law costume.

\section{Remove the discretion first}

The most important thing in this chapter, and the one most often left out of
anti-corruption programmes because it does not look like anti-corruption.

Corruption requires a discretionary decision. Where an official has no choice ---
where the rule is published, the rate is single, the answer is automatic, the
deadline is binding and the appeal is real --- there is nothing to sell. Most
effective anti-corruption is therefore deregulation and simplification rather
than enforcement, and it has the additional advantage of being achievable by a
government that does not yet have honest courts.

Cuba is an extreme case because the number of discretionary surfaces is
extraordinary. Every one of the following has been, and mostly still is, a
decision an official takes about a citizen: whether an activity may be licensed;
which category it falls in; whether a permit is renewed; whether an import is
allowed; what exchange rate applies; whether currency is allocated; whether a
building may be extended; whether a plot is granted; whether a business may be
inspected today; whether a doctor may travel.

\begin{design}{Discretion, and its removal}
\label{d:discretion}
\begin{rules}
\rl{1}{Rules, not permissions} Activity is permitted by right subject to
published prohibitions (Design~\ref{d:commit}.2). Every remaining licensing
regime must be individually justified, published, and reviewed on a sunset
schedule.

\rl{2}{Single rates} One exchange rate (Design~\ref{d:regime}.3), one VAT rate
(Design~\ref{d:tax}.2), one tariff schedule, one set of published fees. Every
multiple rate in Cuban history has become a rent, and the largest single
corruption surface in the economy for thirty years has been access to foreign
currency at the official rate.

\rl{3}{Deemed consent} Statutory decision periods on every authorisation, with
expiry constituting a grant (Design~\ref{d:invest}.2). An official who cannot
delay cannot extract, and this single mechanism does more work than an agency.

\rl{4}{No negotiated terms} No sectoral tax holidays, no bespoke incentive
packages, no individually negotiated rates or exemptions
(Design~\ref{d:tax}.1). Anything negotiated one-to-one is negotiated by two
people in a room.

\rl{5}{Appeal that works} Every refusal, delay, seizure and penalty appealable to
the administrative court of Design~\ref{d:courts}.4, with the state bearing costs
when it loses. This is what converts the rules above from statements into
constraints.

\rl{6}{Sunset} Every regulation carries an expiry date and must be re-enacted.
Regulatory accumulation is how discretion returns after a reform, quietly, one
decree at a time.
\end{rules}
\end{design}

\section{Publish everything}

The second instrument, and the cheapest.

\begin{design}{Transparency by default}
\label{d:transparency}
\begin{rules}
\rl{1}{Open contracting} Every public contract, tender, award and amendment
published in a machine-readable open standard, from day one, with the bidders,
the evaluation and the price. This includes contracts of state enterprises and of
the entities in Design~\ref{d:gaesa}, without exception. Cuba is starting from
zero, which is an advantage: there is no legacy system to migrate.

\rl{2}{Asset declarations} Annual public declarations of income, assets,
liabilities and interests by ministers, deputies, judges, senior officials,
enterprise directors, military officers above a rank, and their household
members. Verified by sampling against the registries, with a published sanction
for false declaration.

\rl{3}{Beneficial ownership register} A public register of the ultimate natural
persons behind every company, trust and vehicle operating in Cuba, with no
anonymous ownership permitted and no exceptions
(Design~\ref{d:privat}.5).

\rl{4}{Budget and accounts} The single budget of Design~\ref{d:fiscal}.1 and its
execution published quarterly, at national and municipal level
(Design~\ref{d:municipal}.4), in a common format.

\rl{5}{Registries open} Cadastre, property register, company register, licences
granted, land allocations, and privatisation results all publicly searchable.

\rl{6}{Priority targets} Where capacity is scarce, publication effort goes first
to the sectors with the largest rents: energy procurement
(Design~\ref{d:energyseq}.4), tourism and hotel investment, port and logistics
concessions, mining, and the divestments under Design~\ref{d:gaesa}.5.
\end{rules}
\end{design}

Publication only works if somebody reads it, which is why
Design~\ref{d:minimum}.2 and \ref{d:minimum}.3 --- a free press and a statutory
right to information --- are in Chapter~\ref{ch:polity} rather than here. An open
contracting portal in a country with no journalists is a database nobody opens.
This is the single most important dependency in the chapter and the reason the
political sequencing cannot be deferred: \emph{the enforcement mechanism for
transparency is publicity, and publicity requires a press.}

\section{Cash, and the digital state}

Corruption's medium is cash and its enabler is the paper transaction that leaves
no record. Cuba is an overwhelmingly cash economy, which is one reason
Design~\ref{d:pay2}'s payment rails appear in the money chapter as
infrastructure rather than convenience.

\begin{design}{Removing the cash}
\label{d:digital}
\begin{rules}
\rl{1}{Public payments digital} All state payments and receipts above a low
threshold --- wages, pensions, benefits, taxes, fees, supplier payments ---
electronic and traceable. This also delivers Design~\ref{d:libreta}'s conversion
and Design~\ref{d:tax}'s withholding.

\rl{2}{Procurement digital} All public purchasing through a single electronic
platform with automatic publication, standard documents, and an audit trail that
cannot be edited.

\rl{3}{Customs} An electronic single window with published tariffs, risk-based
inspection selected by algorithm rather than by officer, and published clearance
times. Customs is the highest-yield corruption surface in every small importing
economy, and Cuba is about to become one.

\rl{4}{Interoperable registries} Tax, property, company, beneficial ownership and
asset declarations queryable against each other by the audit and anti-corruption
bodies. This is what makes clause~2 of Design~\ref{d:transparency} verifiable
rather than ceremonial.
\end{rules}
\end{design}

\section{The institutions, and their honest limits}

The conventional anti-corruption programme --- an agency, a law, a hotline ---
comes last here deliberately, because it is the part that fails most often and
that reformers reach for first.

\begin{design}{Enforcement}
\label{d:enforce}
\begin{rules}
\rl{1}{Auditor-general} An independent supreme audit institution reporting to the
assembly and not to the executive, with its own budget, the power to audit any
public body including the entities in Design~\ref{d:gaesa}, and an obligation to
publish every report in full and on time.

\rl{2}{Anti-corruption agency} An independent body with its own budget, its own
investigators and its own prosecutors, appointed on the model of
Design~\ref{d:courts}.1 and removable only by the process in
Design~\ref{d:courts}.2. \emph{With the caveat stated plainly}: agencies of this
kind fail without independent courts to prosecute into, and where they have been
created ahead of judicial reform they have been captured and used against
political opponents, which is worse than not having one. This clause is
conditional on Design~\ref{d:courts} being in force.

\rl{3}{Whistleblowers} Statutory protection for disclosure in the public
interest, in the public and private sectors, with reinstatement, compensation and
anonymity, and a defence to any charge under the state-secrets provisions.

\rl{4}{Asset recovery} A civil-recovery route allowing the state to recover
unexplained assets on the balance of probabilities, with judicial supervision and
a right of explanation. Per Design~\ref{d:amnesty}.4, recovery of the asset is the
remedy for historic enrichment rather than imprisonment --- more achievable,
more useful, and it does not require the state to relitigate the past.

\rl{5}{Pay the officials} A tax inspector, customs officer, judge or procurement
official who cannot live on their salary will supplement it. Public pay under
Design~\ref{d:pay} is an anti-corruption instrument and is the one most likely to
be cut when money is short, which is exactly when it matters.
\end{rules}
\end{design}

\section{Measurement}

A design that cannot be checked is a wish. Chapter~\ref{ch:sequence} collects the
book's indicators; the ones belonging to this chapter are listed here because
they are the earliest warning available and because they are cheap.

\begin{design}{What gets published annually}
\label{d:measure}
\begin{rules}
\rl{1}{Concentration} The share of privatised assets by value acquired by the
largest ten beneficial owners, and the share acquired by persons who held public
office or a state enterprise directorship in the preceding five years. If this
number is large, the design is failing in the way it is most likely to fail.

\rl{2}{Competition} The share of public contracts by value awarded through open
competitive tender with more than two bidders, and the share awarded directly.

\rl{3}{Discretion} The number of licensing and authorisation regimes in force, and
the share of authorisations decided within the statutory period. Both should fall
and rise respectively; if the count of regimes is rising, discretion is returning
by accumulation.

\rl{4}{Compliance} The share of asset declarations filed on time and the number
sampled, verified and sanctioned.

\rl{5}{Independent measurement} Cuba should join and be assessed by the
international transparency and open-contracting standards, and publish the
assessments whatever they say. External measurement is a commitment device in
the sense of Design~\ref{d:commit}: it is a judgment the government cannot
edit.
\end{rules}
\end{design}

\section{The order}

To close, the ordering claim of the chapter, since order is what it is really
about.

Discretion is removed \emph{before} assets are sold. Accounts are published
\emph{before} the entities holding them are restructured. The press and the right
to information are established \emph{before} the transparency portals, because
the portals are useless without them. The courts come \emph{before} the
anti-corruption agency. And the largest assets are sold \emph{last}, when all of
the above exist.

Reverse any of these and the instrument does not merely fail; it becomes an
instrument of the thing it was meant to prevent --- an anti-corruption agency
that prosecutes opponents, an auction that launders a pre-agreed allocation, a
transparency portal that publishes into silence. That is the specific way this
chapter goes wrong, and it goes wrong through haste rather than through bad
intent.
